\chapter{Топологія: O-точка, X-точка, сепаратриса}\label{ch:geo5}

\section{Критичні точки і гесіан}\label{sec:geo5-crit}

Поверхні $\psi=\text{const}$ «ламаються» лише там, де $\grad\psi=0$, тобто де $\Bpol=0$. Такі точки називають критичними. Їхній тип визначає гесіан $\mathsf{H}_\psi=\bigl(\begin{smallmatrix}\psi_{RR}&\psi_{RZ}\\\psi_{RZ}&\psi_{ZZ}\end{smallmatrix}\bigr)$:
\begin{itemize}
\item $\det\mathsf{H}_\psi>0$ --- \textbf{O-точка} (еліптична): екстремум $\psi$, навколо --- вкладені замкнені контури. Це магнітна вісь або центр магнітного острова.
\item $\det\mathsf{H}_\psi<0$ --- \textbf{X-точка} (гіперболічна, сідло): крізь неї проходять дві гілки лінії рівня, що перетинаються. Ця лінія рівня --- \emph{сепаратриса}.
\end{itemize}
Докладніше (індекс Пуанкаре--Хопфа, теорія Морса) --- \cite[розд.~4, п.~4.4]{Manual}. Скрипт гайду шукає X-точку як мінімум $|\grad\psi|^2$ на сплайні і рахує $\det\mathsf{H}_\psi$ з других похідних сплайна:

\code{geo/geom_tour.py}{72}{85}{(X-точка і гесіан)}

\noindent Рядок 74: цільова функція $|\grad\psi|^2$ має мінімум нуль у \emph{будь-якій} критичній точці, тож стартова точка вирішує, яку з них знайде метод. У рядку 101 файлу старт --- найнижча точка LCFS. Рядки 81--85: $\det\mathsf{H}_\psi$ з аналітичних других похідних бікубічного сплайна. На сітці з кроком 2,7--5~см це оцінка, а не точне значення.

\section{Магнітна вісь зблизька}\label{sec:geo5-axis}

Біля O-точки $\psi-\psiax\approx\frac12\,\vect{x}^{\mathsf T}\mathsf{H}_\psi\vect{x}$, де $\vect x=(R-\Rax,Z-\Zax)$: поверхні --- подібні еліпси. Звідси два «локальні» параметри геометрії.
\begin{itemize}
\item \emph{Витягнутість на осі} $\kappa_0=\sqrt{\psi_{RR}/\psi_{ZZ}}$ (при $\psi_{RZ}\approx0$).
\item \emph{$q$ на осі.} Еліпс рівня $\delta\psi$ має площу $2\pi\,\delta\psi/\sqrt{\det\mathsf{H}_\psi}$. Тороїдальний потік через нього на радіан $\psit=(F/\Rax)\cdot\text{площа}/2\pi$, тому
\begin{equation}\label{eq:geo5-q0}
q_0=\frac{\dd\psit}{\dd\psi}=\frac{F}{\Rax\sqrt{\det\mathsf{H}_\psi}} .
\end{equation}
Цю формулу використовує \texttt{SolovevEquilibrium.q\_axis} для JET (розд.~\ref{ch:geo6}).
\end{itemize}

Для DIII-D, кадр 75: $\psi_{RR}=3{,}47$, $\psi_{ZZ}=2{,}29$, $\psi_{RZ}=-0{,}04$~Вб/(рад$\cdot$м$^2$), $\det\mathsf{H}_\psi=7{,}95>0$. Отже, $\kappa_0=1{,}23$ проти $\kappa=1{,}70$ на межі: \emph{витягнутість зростає від осі до краю}. Це типово, бо форму задають зовнішні котушки, і їхній вплив згасає вглиб плазми. За~\eqref{eq:geo5-q0} $q_0=0{,}68$, тоді як контурний інтеграл на $\psiN=0{,}02$ дає 0,711. Різниця $\approx4\,\%$ --- це межа точності других похідних сплайна на сітці $65\times65$ біля осі, а не фізика.

\section{X-точка і сепаратриса}\label{sec:geo5-x}

Для кадру 75 скрипт знаходить X-точку в $(R,Z)=(\SI{1,2407}{м},\,\SI{-1,0962}{м})$. Там $\Bpol=1{,}6\cdot10^{-9}$~Тл (чисельний нуль) і $\psiN=1{,}00002$: X-точка лежить \emph{на} LCFS і збігається з її найнижчою точкою (рис.~\ref{fig:geo4-diiid}). Це однонульова діверторна конфігурація (lower single null): межу плазми задає не стінка, а сепаратриса.

Гесіан у X-точці: $\psi_{RR}=0{,}463$, $\psi_{ZZ}=-0{,}497$, $\psi_{RZ}=-0{,}546$; власні значення $-0{,}743$ і $+0{,}710$, $\det\mathsf{H}_\psi=-0{,}53<0$ --- сідло. Звідси маленька, але гарна перевірка фізики.

\begin{derivation}[чому гілки сепаратриси перетинаються під прямим кутом]
У X-точці $\psi_R=0$, тому оператор Ґреда--Шафранова $\GSop\psi=\psi_{RR}-\psi_R/R+\psi_{ZZ}$ зводиться до сліду гесіана, а рівняння рівноваги~\cite[розд.~3]{Manual} дає $\psi_{RR}+\psi_{ZZ}=-\mu_0R_X\Jtor$, де $R_X$ --- радіус X-точки. На сепаратрисі струму майже немає, тож слід малий, і власні значення протилежні: $\lambda_1\approx-\lambda_2$. Нульова лінія квадратичної форми $\lambda_1u^2+\lambda_2v^2=0$ --- це пара прямих $v=\pm\sqrt{-\lambda_1/\lambda_2}\,u$, і при $\lambda_1=-\lambda_2$ вони перетинаються під кутом $90^\circ$. Для кадру 75 слід дорівнює $-0{,}034$ (5\,\% від $|\lambda|$), кут між гілками $\approx91^\circ$, а $|\Jtor|\approx0{,}034/(\mu_0\cdot\SI{1,24}{м})\approx\SI{2e4}{А/м^2}$. Для порівняння, середня густина струму $\Ip/S=\SI{1,36}{МА}/\SI{1,77}{м^2}\approx\SI{7,7e5}{А/м^2}$ ($|\Ip|=\SI{1,359}{МА}$ --- за законом Ампера, циркуляцією $\Bpol$ уздовж LCFS, функція \texttt{ip\_ampere} у \texttt{geom\_tour.py}, рядки 64--69), тобто в X-точці лише $\sim3\,\%$ від середньої. Точність других похідних тут та сама, що й на осі, тому число --- порядок величини.
\end{derivation}

Навколо сепаратриси простір поділено на три області:
\begin{itemize}
\item \emph{всередині LCFS} --- замкнені вкладені поверхні, «ядро» плазми;
\item \emph{SOL} (scrape-off layer) --- поверхні $\psiN>1$ над X-точкою, відкриті: силова лінія йде вздовж плазми й упирається в пластини дивертора;
\item \emph{приватна область} (private flux region) --- під X-точкою, між двома «ногами» сепаратриси. Тут також $\psiN<1$, хоча плазми ядра там немає. Тому маска «$\psiN<1$» без перевірки «всередині LCFS» захоплює цю область; на рис.~\ref{fig:geo3-thetastar} її довелося вирізати окремо.
\end{itemize}

\emph{Чому $q\to\infty$ на сепаратрисі.} У~\eqref{eq:geo3-q} підінтегральна функція $1/(R|\grad\psi|)$ біля X-точки поводиться як $1/(\text{відстань до X})$, і для поверхні, що проходить на відстані $d$ від X-точки, інтеграл росте як $\ln(1/d)$. Лінія, яка проходить поблизу X-точки, «застрягає» там на багато тороїдальних обертів. Звідси крутий ріст $q$ біля межі: $q(0{,}95)=3{,}70$, $q(0{,}98)=4{,}50$ (розд.~\ref{ch:geo3}).

\section{Як знаходять LCFS}\label{sec:geo5-lcfs}

EFIT видає LCFS разом з рівновагою, але оцінювач конкурсу Fusion Equilibrium Challenge витягає її сам із будь-якої карти $\psi$ --- і справжньої, і передбаченої. Означення: LCFS --- найнижчий рівень $\texttt{AXIS\_SIGN}\cdot\psi$, який ще дає \emph{один замкнений контур навколо осі всередині допустимої області машини}. Нижче цього рівня контур або виходить на стінку (лімітер), або розпадається в X-точці (дивертор). Рівень шукають бісекцією:

\code{fec/starter/fusion_scoring/lcfs.py}{58}{89}{(бісекція рівня LCFS)}

\noindent Рядки 68--70: вісь --- найсильніший \emph{локальний} максимум, а не глобальний (коментар у рядках 64--67: на 1,4\,\% кадрів MAST поле котушок на краю області сильніше, ніж на осі). Рядки 71--73: нижня межа бісекції --- мінімум $\phi$ на краю сітки (там контур точно відкритий), верхня --- трохи нижче осі (там точно малий замкнений контур). Рядки 80--87: 22 кроки бісекції звужують інтервал у $2^{22}\approx4\cdot10^6$ разів. \texttt{\_best\_contour\_at} (рядки 44--55) приймає контур, лише якщо він замкнений, охоплює вісь і лежить у масці камери. Цей алгоритм однаково працює для лімітерної і діверторної плазми, бо перевіряє саме топологію (замкненість і охоплення), а не тип межі.

\section{Скелет критичних точок як критерій якості}\label{sec:geo5-skeleton}

Справжня рівновага всередині LCFS має рівно одну O-точку і жодної X-точки. Шум у передбаченій $\psi$ народжує зайві пари O--X, тобто острови, яких немає. Проєкт використав це як незалежний критерій якості передбачення; повний розбір із числами (E9--E16, R5, R7, R9) є в методичці~\cite[розд.~4, п.~4.4]{Manual} і путівнику~\cite{Guide}. Тут --- лише код, що рахує індекс:

\code{ourtry/04-novelty/topology_probe.py}{63}{90}{(критичні точки через число обертів $\grad\psi$)}

\noindent Рядки 65--66: напрямок $\grad\psi$ у кожному вузлі. Рядки 69--78: для кожного вузла --- число обертів цього напрямку по 8 сусідах ($+1$ для O-точки, $-1$ для X-точки, 0 для звичайної точки); стрибки кута зводяться в $(-\pi,\pi]$ (рядок 77). Рядки 80--89: одна справжня критична точка «засвічує» блок $\approx2\times2$ вузли, тому сусідні вузли об'єднують у компоненту і рахують число обертів по периметру її обмежувального прямокутника. Так кожну точку враховано один раз.

\begin{established}{E9--E11}
На кадрах DIII-D \#203702 без шуму скелет --- 1 O-точка, 0 X-точок, сума індексів $+1$. Скелет стійкий до шуму 3\,\% від $\mathrm{std}(\psi)$. При 10\,\% з'являється $\approx20$ зайвих точок, але сума індексів лишається $\approx+1$, бо точки народжуються парами O--X. Тому інформативна \emph{кількість} точок, а не їхня сума~\cite{RegEst}, \cite[розд.~4]{Manual}.
\end{established}

\begin{practicum}[топологія]
(1)~Запустіть \texttt{find\_x\_point} зі старту на вісі. Що знайде метод і чому? (2)~Знайдіть другу (верхню) критичну точку DIII-D поза LCFS, стартувавши над верхньою точкою межі. Яке в неї $\psiN$ і $\det\mathsf{H}_\psi$? Наскільки розряд далекий від double null? (3)~У \texttt{topology\_probe.py} вимкніть ерозію (\texttt{erode=0}) і поясніть, чому сума індексів стає $1-k_X$.
\end{practicum}
